\subsection{So sánh sản phẩm với các giải pháp hiện có}

\subsubsection{So sánh nhóm đối thủ trực tiếp}

Bảng dưới đây chỉ so sánh iLive, AnyLive và 1Stream vì đây là nhóm có phạm vi
workflow AI livestream tương đối gần nhau. Ký hiệu không được hiểu là xếp hạng
chất lượng:
\begin{itemize}
    \item \textbf{Đã thương mại hóa}: năng lực được công bố và đang cung cấp;
    \item \textbf{Có một phần}: có hỗ trợ nhưng không phải trọng tâm hoặc cần cấu hình;
    \item \textbf{MVP dự kiến}: nằm trong phạm vi xây dựng của 1Stream nhưng chưa qua paid pilot;
    \item \textbf{Ngoài MVP}: chưa triển khai trong giai đoạn hiện tại.
\end{itemize}

\begin{table}[H]
    \centering
    \footnotesize
    \setlength{\tabcolsep}{4pt}
    \renewcommand{\arraystretch}{1.28}
    \caption{So sánh nhóm đối thủ AI livestream trực tiếp}
    \label{tab:direct-competitor-comparison}
    \rowcolors{2}{paperBg}{white}
    \begin{tabularx}{\linewidth}{@{}L{4.2cm}L{3.0cm}L{3.0cm}Y@{}}
        \toprule
        \rowcolor{softBlue}
        \textbf{Tiêu chí} & \textbf{iLive} & \textbf{AnyLive} &
        \textbf{1Stream MVP} \\ \midrule

        AI avatar/video livestream &
        Đã thương mại hóa; có kho host và custom host &
        Đã thương mại hóa; AI live streamer 24/7 &
        MVP dự kiến; dùng avatar/giọng có sẵn \\

        Tạo và cập nhật kịch bản &
        Có kịch bản theo gói, FAQ/keyword &
        Có tạo và cải tiến script dựa trên dữ liệu &
        MVP dự kiến từ dữ liệu khóa học đã duyệt \\

        Phát dài giờ/đa nền tảng &
        Có gói 30--90 giờ; một số gói đa kênh &
        Có 24/7, đa nền tảng và đa ngôn ngữ &
        Một kênh trong MVP; chưa cam kết 24/7 \\

        Phản hồi bình luận realtime &
        Có FAQ/keyword theo phạm vi gói &
        Có phản hồi và phân tích comment &
        FAQ cơ bản dự kiến; câu khó chuyển người thật \\

        Analytics và cải tiến sau phiên &
        Có báo cáo/hỗ trợ theo gói; mức chi tiết tùy gói &
        Năng lực cốt lõi: sales, viewership, comments &
        Báo cáo cơ bản dự kiến; chưa có dữ liệu thực tế \\

        Dịch vụ triển khai và vận hành &
        Đã có quy trình, hỗ trợ và chính sách bù giờ &
        Có nền tảng và tư vấn triển khai cho thương hiệu &
        Vận hành bán thủ công trong pilot \\

        Phù hợp ICP trung tâm khóa học &
        Không công bố là phân khúc ưu tiên &
        Không công bố là phân khúc ưu tiên &
        Là phân khúc mục tiêu dự kiến; chưa kiểm chứng \\

        Bằng chứng thị trường &
        Có gói giá công khai và case thương hiệu &
        Có sản phẩm/case ở nhiều thị trường &
        Chưa có paid pilot hoàn tất \\ \bottomrule
    \end{tabularx}
    \rowcolors{2}{}{}
\end{table}

Nhận xét chính:
\begin{itemize}
    \item iLive và AnyLive hiện mạnh hơn 1Stream về độ trưởng thành, số giờ phát,
    đa nền tảng và bằng chứng triển khai;
    \item 1Stream chỉ có giả thuyết khác biệt ở ICP hẹp, dữ liệu khóa học, gói vào
    nhỏ và flow chuyển lead tuyển sinh;
    \item vì chưa có paid pilot, cột 1Stream không được diễn giải là ``Có'' theo
    nghĩa năng lực đã được chứng minh.
\end{itemize}

\subsubsection{Phân tích giải pháp thay thế theo cấu phần}

Các giải pháp dưới đây không được chấm bằng bộ tiêu chí của nhóm AI livestream.
Thay vào đó, mỗi giải pháp được đánh giá theo công việc mà khách hàng có thể dùng
nó để thay thế.

\begin{table}[H]
    \centering
    \footnotesize
    \setlength{\tabcolsep}{4pt}
    \renewcommand{\arraystretch}{1.25}
    \caption{Các giải pháp thay thế theo từng cấu phần}
    \label{tab:component-substitutes}
    \rowcolors{2}{paperBg}{white}
    \begin{tabularx}{\linewidth}{@{}L{2.8cm}L{3.2cm}Y Y@{}}
        \toprule
        \rowcolor{softBlue}
        \textbf{Giải pháp} & \textbf{Công việc thay thế} &
        \textbf{Điểm mạnh trong phạm vi riêng} &
        \textbf{Phần khách hàng vẫn phải tự làm} \\ \midrule

        GoStream &
        Phát/lên lịch từ video có sẵn &
        Tập trung vào streaming, dễ dùng khi nội dung đã chuẩn bị xong &
        Chuẩn hóa dữ liệu, tạo video, duyệt, FAQ, thu lead và báo cáo kinh doanh \\

        Botcake &
        Chatbot, FAQ, thu lead và handover &
        Messaging flow, lead automation và chuyển bot--nhân viên &
        Tạo video, phát livestream và giám sát giờ-kênh \\ \bottomrule
    \end{tabularx}
    \rowcolors{2}{}{}
\end{table}

\subsubsection{Phân tích giải pháp gián tiếp}

\begin{table}[H]
    \centering
    \footnotesize
    \setlength{\tabcolsep}{4pt}
    \renewcommand{\arraystretch}{1.25}
    \caption{Các giải pháp gián tiếp trong social commerce và quản lý bán hàng}
    \label{tab:indirect-solutions}
    \rowcolors{2}{paperBg}{white}
    \begin{tabularx}{\linewidth}{@{}L{2.8cm}L{3.3cm}Y Y@{}}
        \toprule
        \rowcolor{softBlue}
        \textbf{Giải pháp} & \textbf{Phạm vi chính} &
        \textbf{Mức liên quan với 1Stream} &
        \textbf{Lý do không đặt trong bảng đối thủ trực tiếp} \\ \midrule

        Harasocial &
        Hội thoại, bình luận, livestream commerce và tạo đơn &
        Có thể thay thế một phần khâu phản hồi và quản lý tương tác &
        Không tạo/vận hành video AI theo giờ-kênh; trọng tâm là commerce \\

        Sapo &
        Chat AI, đơn hàng, tồn kho và quản lý bán hàng hợp kênh &
        Có thể nhận lead hoặc tiếp tục chăm sóc sau phiên &
        Phạm vi là hệ thống bán hàng rộng; không phải AI livestream service \\ \bottomrule
    \end{tabularx}
    \rowcolors{2}{}{}
\end{table}

\subsubsection{Kết luận phương pháp so sánh}

Không có một bảng duy nhất có thể phản ánh công bằng tất cả giải pháp:
\begin{itemize}
    \item iLive, AnyLive và 1Stream được so sánh theo workflow AI livestream;
    \item GoStream và Botcake được phân tích theo mức độ thay thế từng cấu phần;
    \item Harasocial và Sapo được xem là giải pháp gián tiếp, không chấm cùng thang;
    \item agency và đội ngũ nội bộ được xem là phương án thay thế về vận hành.
\end{itemize}

Cách trình bày này tránh thiên kiến ``1Stream có nhiều ô xanh nhất'' và giúp chỉ
ra rõ hơn rủi ro thật: khách hàng có thể chọn một giải pháp trọn bộ trưởng thành,
hoặc tự ghép các công cụ chuyên biệt với chi phí và công vận hành khác nhau.
